
Layouters are used to arrange the nodes and the edges of a graph
on the screen to optimally suit the representation of the
data the graph contains. As there are a huge number of
different data a graph may contain and an even larger number of
opinions on what an optimal representation of this data is, there
exists a huge amount of different layout algorithms. In the
ffgraph library layout algorithms are implemented as graph
labels. By this you can add specific layouters to specific graphs
and other classes like windows can determine at runtime which
layouters are available. These classes can even use layouters
they did not know about when they were written. E.g. the
graph-window class {\bf FFTclWindow} included in the distribution
automatically adds all layouters available for a graph to its
layout menu.

\subsection{Using layouters} 
\label{using_layouters}

Layouters are usually subclasses of {\bf FFLayoutLabel} which is
a subclass of {\bf FFGraphLabel}. Of course layouters do not need
to be implemented as graph labels.  This subsection only
describes how to use the layouters included in the library. It
should work for other layouters as well, but may not.

In your program you have to add layout labels to all graphs you
want to layout using the {\bf
  FFGraphClass::addGraphLabel(FFGraphLabel *)} method.  E.g.

\bigskip

{\tt  
FFGraph g;\\
g->addGraphLabel(new FFSugi);\\
g->addGraphLabel(new FFSpring);\\
}

\bigskip

You can then use the layouter using one of the following methods:

\begin{itemize}

\item{\bf void FFLayoutLabel::layout()}

The whole graph is layouted.

\item{\bf void FFLayoutLabel::layout(AVLSet\verb|<|FFNode\verb|>|\& set)}

Only the subgraph defined by the nodes in {\bf set} is layouted.

\end{itemize}

If you want to change the parameters of a layouter, you can
either use the methods the layouters provide\footnote{see
  sections \vref{ffsugi} and \vref{ffspring}}, or if you use {\bf
  FFTclWindow} to display the graph, you can set the parameters
at runtime using the inspector included in the library.

If you use the class {\bf FFTclWindow} to display the graph, the
new layouter is automatically added to its layout menu.

\subsection{Writing layouters}
\label{writing_layouters}

The easiest way to write your own layouter is to subclass {\bf
  FFLayoutLabel} and implement these methods:

\begin{itemize}

\item{\bf void FFLayoutLabel::layout()}

To layout the whole graph.

\item{\bf void
    FFLayoutLabel::layout(AVLSet\verb|<|FFNode\verb|>|\& set)}

To layout the subgraph defined by the nodes in {\bf set}.

\end{itemize}

If you use {\bf FFTclWindow} to display the graph, you only have
to add your layouter to the graph using {\bf int
  FFGraphClass::addGraphLabel(FFGraphLabel *)} and it will
automatically be added to the layout menu of {\bf FFTclWindow}. If the
menu is invoked, the method {\bf void FFLayoutLabel::layout()}
is called. If more than one node is selected, {\bf void
FFLayoutLabel::layout(AVLSet\verb|<|FFNode\verb|>|)} is called.

That is all you basically have to do.  If you want the user to be
able to set parameters at runtime the usual way to do this is
to use an inspector. To use an inspector you have to overwrite
the method

\begin{itemize}
\item{\bf BOOL FFLabelClass::isInspectable()}
\end{itemize}

to return {\bf YES}.  

You then have to decide if you want to use your own inspector or
if you want to use the inspector included in the library.  If you
want to use the given inspector, you just have to implement the
methods:

\begin{itemize}

\item {\bf string FFLabelClass::unparse()}\\
  and \\
  {\bf BOOL FFLabelClass::parse(String\&)}

to read and write the layout parameters in string form.  
 
\end{itemize}

If you want to use your own inspector, you have to overwrite
the method 

\begin{itemize}
\item{\bf BOOL FFGraphLabel::hasOwnInspector()}
\end{itemize}

to return {\bf YES}. You also have to write your own inspector
class. You can do this  subclassing {\bf FFInspector}\footnote{See
section \vref{writing_inspectors} for further details}.
